\subsection{Shared Canonical Note-Wise Comparison}

We compare Polytune~\citep{chou2025polytune}, prompted
LadderSym~\citep{chou2026laddersym}, and our method on the same complete
358-clip outputRaw validation set. The frozen manifest and canonical target
records define identical sample membership and ground truth for all three
methods. All evaluation performance WAVs and MIDIs were verified against the
manifest's SHA-256 hashes. Of these inputs, 26 were available locally and
332 were reconstructed from their original source scores and sample seeds;
all reconstructed audio and MIDI files matched byte for byte, and the
clean/performed/deleted note lineage matched the audited targets.

\paragraph{Training and checkpoint selection.}
Our result uses the previously reported exact-32-dimensional checkpoint.
The audited outputRaw release provides 4,544 training records. We evaluate the
retained best Polytune and prompted LadderSym checkpoints after 31 and 26
completed epochs, respectively. Both baselines were initialized from scratch
and trained on a separate 10,049-clip synthetic corpus, with checkpoint
selection on its 1,116-clip validation split. Their configuration uses seed
365, a peak learning rate of $2\times10^{-5}$, 4,000 warmup updates, a cosine
schedule originally configured for 40 epochs, and 16-step gradient
accumulation with one sampled 2.048-second segment per batch. Inference uses
16\,kHz mono inputs, float32 greedy decoding, and a 1,024-token output budget.
LadderSym additionally uses a 1,024-token reference-score prompt. The complete
recordings are evaluated.

These existing checkpoints have different training exposure. An exact-audio
hash audit found that 21 of the 358 evaluation recordings occurred in baseline
training and three occurred in its original checkpoint-selection set. We
retain the full common evaluation list and disclose these overlaps; the
comparison describes the available trained systems and does not isolate
architecture under a matched training budget or wholly unseen evaluation
recordings. All clips are Mozart snippets, so this evaluation also does not
establish generalization to unseen works.

\paragraph{Shared score-location conversion.}
The two baselines predict timestamped Extra, Missing, and Correct notes.
We apply the same fixed conversion to both outputs before scoring. Mutually
nearest Extra and Missing notes with different pitches and onsets within
100\,ms form a substituted-note prediction. Correct notes, these substitutions,
and remaining Missing notes are aligned monotonically to canonical events in
the clean reference score. The edit alignment uses unit insertion/deletion
costs and a pitch-mismatch cost of 1.5; a substitution uses the Missing note's
expected pitch as its alignment anchor. Gaps in this conversion do not create
additional Missing predictions.

Remaining Extra runs of at least three notes can identify a replay of a
previously visited contiguous score sequence with at least 80\% pitch
agreement. The longest qualifying sequence is chosen, with ties resolved by
pitch agreement and then the most recent score position; prior predicted
visits determine the copy pass. The conversion uses only model predictions
and the clean reference score. Its parameters are fixed without validation-score
tuning. A process-level guard prevents access to labels and target databases
during conversion, and predictions are frozen and hash-verified before a
separate scoring process opens the gold targets. Unlocated and duplicate
predictions remain in the prediction count. In particular, no gold rendered-note
identity is assigned to an unlinked Extra prediction.

\paragraph{Unified metric.}
All three methods in Figure~\ref{fig:canonical-validation} use the implemented
canonical combined-pipeline metric. The location is the canonical score event
or event span together with its copy pass, or an audited rendered-event
identity for an extra note. The complete event type comprises the alignment
relationship, note-error class, rhythm flag, and replay flag; deleted score
events are scored as missed notes. Rhythm remains part of this common target,
while the baselines emit no rhythm-error detections. Intonation is masked for
all methods.

The same exclusive one-to-one matcher assigns credit 1 for the correct location
and complete type, 0.5 for the correct location with a different type, and 0
for an incorrect or unavailable location. We pool credit, prediction counts,
and the same 40,712 target events over all 358 clips to obtain micro precision,
recall, and F1. Confidence intervals use clip bootstrap resampling; the new
baseline runs use 1,000 replicates with seed 365. Unit checks verify that the
baseline scoring path agrees with our full-pipeline scorer for the same
canonical labels. We report this single common metric for the comparison.

\begin{figure}[htbp]
\centering
\includegraphics[width=\linewidth]{figures/canonical-validation.pdf}
\caption{Canonical combined-pipeline micro-F1 on the same 358 outputRaw
validation clips. Whiskers show 95\% clip-bootstrap intervals.
Both baselines include the shared score-location conversion. Training corpora
differ, and 21 evaluation recordings occurred in baseline training.
The comparison therefore describes the available checkpoints under the
training exposure detailed in the text.}
\label{fig:canonical-validation}
\end{figure}

Our result retains the previously completed canonical validation measurement;
the two baselines are newly evaluated using their fixed checkpoints. The comparison
includes errors introduced by each system's score localization. It should be
interpreted with the distinct training corpora and baseline training overlap
stated above.
